%% ---------------------------------------------------------------------
%%  Le Potenze: Schema Visivo — le regole delle operazioni e i casi
%%  particolari in UNA pagina A4.
%%
%%  Ordine dell'operatore, 27/11/2026: «stesso lavoro su questo PDF: le
%%  proprietà delle potenze, stessa coerenza e stessa pubblicazione». Il PDF
%%  è il suo («Le Potenze: Schema Visivo», 110034 B): le due sezioni e le
%%  otto proprietà sono quelle sue, nessuna tolta per far stare la pagina —
%%  il respiro lo toglie la libreria (\schedaCompatta, \schedaMarginiSicuri).
%%
%%  LA LEGENDA DEL SUO PDF, TRADOTTA NELLE FORME DELLA CASA:
%%      base 1 «a»  ->  \formaQuadrato   (blu)
%%      base 2 «b»  ->  \formaTriangolo  (rosso)
%%      esponenti «n», «m»  ->  restano LETTERE: sono esponenti, non basi,
%%      e una forma in apice non si legge (la mappatura è nel report).
%%
%%  Qui non c'è grafica: tutto sta in ../schede.sty.
%% ---------------------------------------------------------------------
\documentclass[10pt]{article}
\usepackage{schede}

%% Testata e piede dentro la zona che ogni stampante sa stampare.
\schedaMarginiSicuri
%% Due sezioni e otto righe in una pagina: si toglie respiro agli SPAZI,
%% non al testo (la regola della libreria).
\schedaCompatta

\schedaTitolo{Le Potenze: Schema Visivo}
             {Operazioni \textperiodcentered\ Esponenti \textperiodcentered\ Casi Particolari}
\schedaPiede{Schema di Matematica — Le Potenze}

\begin{document}

%% ---------------------------------------------------------------------
%%  LA LEGENDA — le forme al posto delle lettere, come nel suo PDF
%% ---------------------------------------------------------------------
{\raggedleft\footnotesize\color{schedeGrigio}%
 Legenda:\quad \formaQuadrato{} base $a$ \quad \formaTriangolo{} base $b$
 \quad $n$, $m$ esponenti\par}
\vspace{2pt}

%% =====================================================================
\schedaSezione{1}{Operazioni con la Stessa Base}
%% =====================================================================
\schedaColonne{0.275}{0.365}{0.36}
\begin{schedaTabella}{Proprietà}{Schema visivo}{Esempio numerico}

  \riga{\voce{Prodotto stesse basi}{(si sommano gli esponenti)}}
       {$\formaQuadrato^{n}\cdot\formaQuadrato^{m}
         \;\freccia\;\formaQuadrato^{n+m}$}
       {$3^{2}\cdot 3^{4}=3^{2+4}=3^{6}$}

  \riga{\voce{Quoziente stesse basi}{(si sottraggono gli esponenti)}}
       {$\formaQuadrato^{n}:\formaQuadrato^{m}
         \;\freccia\;\formaQuadrato^{n-m}$}
       {$5^{7}:5^{3}=5^{7-3}=5^{4}$}

  \riga{\voce{Potenza di potenza}{(si moltiplicano gli esponenti)}}
       {$(\formaQuadrato^{n})^{m}
         \;\freccia\;\formaQuadrato^{n\cdot m}$}
       {$(2^{3})^{2}=2^{3\cdot 2}=2^{6}$}

\end{schedaTabella}

%% =====================================================================
\schedaSezione{2}{Stesso Esponente \& Casi Particolari}
%% =====================================================================
\schedaColonne{0.275}{0.365}{0.36}
\begin{schedaTabella}{Proprietà}{Schema visivo}{Esempio numerico}

  \riga{\voce{Prodotto stessi esponenti}{(si moltiplicano le basi)}}
       {$\formaQuadrato^{n}\cdot\formaTriangolo^{n}
         \;\freccia\;(\formaQuadrato\cdot\formaTriangolo)^{n}$}
       {$4^{2}\cdot 3^{2}=(4\cdot 3)^{2}=12^{2}$}

  \riga{\voce{Quoziente stessi esponenti}{(si dividono le basi)}}
       {$\formaQuadrato^{n}:\formaTriangolo^{n}
         \;\freccia\;(\formaQuadrato:\formaTriangolo)^{n}$}
       {$10^{3}:5^{3}=(10:5)^{3}=2^{3}$}

  \riga{\voce{Esponente zero}{(con base diversa da zero)}}
       {$\formaQuadrato^{0}\;\freccia\;1$}
       {$15^{0}=1$}

  \riga{\voce{Esponente uno}{(la base resta invariata)}}
       {$\formaQuadrato^{1}\;\freccia\;\formaQuadrato$}
       {$8^{1}=8$}

  \riga{\voce{Esponente negativo}{(si inverte la base)}}
       {$\formaQuadrato^{-n}\;\freccia\;\dfrac{1}{\formaQuadrato^{n}}$}
       {$2^{-3}=\dfrac{1}{2^{3}}=\dfrac{1}{8}$}

\end{schedaTabella}

\end{document}
